\documentclass[aspectratio=169,10pt]{beamer}
\usetheme[numbering=fraction,progressbar=frametitle]{metropolis}
\usepackage{fontspec}
\setsansfont{Latin Modern Sans}
\usepackage{amsmath,amssymb,mathtools,graphicx,booktabs}
\usepackage{polyglossia}
\setdefaultlanguage{french}
\setbeamertemplate{navigation symbols}{}
\setbeamerfont{frametitle}{size=\large}
\setbeamerfont{normal text}{size=\normalsize}
\setlength{\parskip}{.45em}
\hypersetup{pdfauthor={},pdfsubject={Lecture slides},pdfkeywords={},pdftitle={Dérivation, variations et optimisation}}
\providecommand{\N}{\mathbb{N}}
\providecommand{\Z}{\mathbb{Z}}
\providecommand{\D}{\mathbb{D}}
\providecommand{\Q}{\mathbb{Q}}
\providecommand{\R}{\mathbb{R}}
\providecommand{\C}{\mathbb{C}}
\providecommand{\vect}[1]{\overrightarrow{#1}}
\providecommand{\vi}{\vec{\imath}}
\providecommand{\vj}{\vec{\jmath}}
\providecommand{\vk}{\vec{k}}
\providecommand{\norm}[1]{\left\lVert #1\right\rVert}
\providecommand{\vc}[2]{\begin{pmatrix}#1\\#2\end{pmatrix}}
\providecommand{\deter}[2]{\det\left(#1,#2\right)}
\providecommand{\abs}[1]{\left\lvert #1\right\rvert}
\providecommand{\dd}{\,\mathrm{d}}
\providecommand{\e}{\mathrm{e}}
\providecommand{\Card}{\operatorname{Card}}
\providecommand{\E}{\operatorname{E}}
\providecommand{\V}{\operatorname{V}}
\providecommand{\degre}{^{\circ}}
\providecommand{\un}[1]{\,\mathrm{#1}}
\providecommand{\sci}[2]{#1\times10^{#2}}
\providecommand{\pgcd}{\operatorname{PGCD}}
\providecommand{\ppcm}{\operatorname{PPCM}}
\title{Dérivation, variations et optimisation}
\subtitle{Cours et exemples}
\date{}
\author{}
\begin{document}
\begin{frame}[plain]
\vfill
{\usebeamerfont{title}\inserttitle\par}
\medskip
{\usebeamerfont{subtitle}\insertsubtitle\par}
\medskip
{\usebeamercolor[fg]{progress bar}\rule{\textwidth}{.5pt}}
\vfill
\end{frame}
\begin{frame}{Le parcours du chapitre}
\textbf{Objectifs.} Calculer une dérivée sur son domaine ; étudier variations et extrema ; comparer des courbes ; optimiser ; exploiter les symétries.

\textbf{Prérequis.} Nombre dérivé, calcul algébrique, signe d’un trinôme, intervalles.

\textbf{Repère annuel.} 1-10 décembre : 8 h, consolidation les 15 et 17
\end{frame}
\begin{frame}{Fonction dérivée et fonctions usuelles}
Une fonction est dérivable sur un intervalle lorsqu’elle l’est en chacun de ses points ; sa dérivée associe à x le nombre \(f\prime(x)\). Les dérivées usuelles sont : constante \(\mapsto0\) ; \(x\mapsto1\) ; \(x^2\mapsto2x\) ; \(x^3\mapsto3x^2\) ; \(\dfrac{1}{x}\mapsto -\dfrac{1}{x^2}\) pour \(x\ne0\) ; \(\sqrt x\mapsto \dfrac{1}{2\sqrt x}\) pour \(x>0\). Pour \(n\in\Z\), \((x^n)\prime=nx^{n-1}\) sur son domaine de dérivabilité ; pour \(n=0\), la dérivée est 0 directement.
\end{frame}
\begin{frame}{Fonction dérivée et fonctions usuelles}
Pour \(f(x)=x^{-3}\), le domaine est \(\R\setminus\{0\}\) et \(f\prime (x)=-3x^{-4}=-\dfrac{3}{x^4}\). Pour \(g(x)=\sqrt x\), la fonction existe en 0, sa dérivée usuelle seulement pour \(x>0\).
\end{frame}
\begin{frame}{Fonction dérivée et fonctions usuelles}
Le domaine de dérivabilité peut être plus petit que le domaine de définition.
\end{frame}
\begin{frame}{Fonction dérivée et fonctions usuelles}
Dériver \(x^5\), \(x^{-2}\) et \(\sqrt x\) en précisant les domaines.
\end{frame}
\begin{frame}{Somme, produit, inverse et quotient}
Pour u et v dérivables, \((u+v)\prime=u\prime+v\prime\), \((\lambda u)\prime=\lambda u\prime\), \((uv)\prime=u\prime v+uv\prime\). Là où \(v\ne0\), \((\dfrac{1}{v})\prime =-\dfrac{v\prime}{v^2}\) et \((\dfrac{u}{v})\prime =\dfrac{u\prime v-uv\prime}{v^2}\).
\end{frame}
\begin{frame}{Somme, produit, inverse et quotient}
Pour \(f(x)=\dfrac{x^2+1}{x-1}\), \(x\ne1\), \(f\prime (x)=\dfrac{[2x(x-1)-(x^2+1)]}{(x-1)^2}=\dfrac{x^2-2x-1}{(x-1)^2}\). Garder les parenthèses évite de changer un seul signe au lieu de tout le numérateur.
\end{frame}
\begin{frame}{Somme, produit, inverse et quotient}
La dérivée d’un produit n’est pas le produit des dérivées ; la dérivée d’un quotient n’est pas leur quotient.
\end{frame}
\begin{frame}{Somme, produit, inverse et quotient}
Dériver \((x^2+1)(3x-2)\) puis \(\dfrac{1}{x^2+1}\).
\end{frame}
\begin{frame}{Signe de la dérivée et extrema}
Sur un intervalle I, si \(f\prime\geq0\), f est croissante ; si \(f\prime\leq0\), elle est décroissante ; si \(f\prime=0\) partout, elle est constante. Une dérivée strictement positive entraîne une stricte croissance. Si f possède un extremum local en un point intérieur où elle est dérivable, sa dérivée y est nulle. La réciproque est fausse.
\end{frame}
\begin{frame}{Signe de la dérivée et extrema}
\begin{columns}[T,onlytextwidth]
\begin{column}{.57\textwidth}
\(f(x)=x^3-3x\) : \(f\prime(x)=3(x-1)(x+1)\). f croît sur \(]-\infty;-1]\), décroît sur \([-1;1]\), croît sur \([1;+\infty[\) ; maximum local \(f(-1)=2\), minimum local \(f(1)=-2\). Aucun de ces extrema n’est global sur \(\R\).
\end{column}
\begin{column}{.4\textwidth}
\centering
\includegraphics[width=\linewidth,height=4.5cm,keepaspectratio]{figures/premiere-06-1.pdf}
\end{column}
\end{columns}
\end{frame}
\begin{frame}{Signe de la dérivée et extrema}
Le théorème s’applique sur un intervalle. Une fonction peut décroître sur deux intervalles séparés sans décroître sur leur réunion.
\end{frame}
\begin{frame}{Signe de la dérivée et extrema}
Étudier \(f(x)=x^2-4x+1\) sur \([0;3]\) et donner ses extrema sur cet intervalle.
\end{frame}
\begin{frame}{Parité, inégalités et comparaison}
Sur un domaine symétrique par rapport à 0, f est paire si \(f(-x)=f(x)\) et impaire si \(f(-x)=-f(x)\). La courbe d’une fonction paire est symétrique par rapport à l’axe des ordonnées ; celle d’une impaire par rapport à l’origine. Pour comparer f et g, on étudie le signe de \(f-g\), éventuellement avec sa dérivée.
\end{frame}
\begin{frame}{Parité, inégalités et comparaison}
Pour \(x>0\), \(h(x)=x+\dfrac{1}{x}-2\) a pour dérivée \(\dfrac{x^2-1}{x^2}\). Elle décroît jusqu’à 1 puis croît ; \(h(1)=0\), donc \(x+\dfrac{1}{x}\geq 2\). L’égalité est obtenue seulement en 1.
\end{frame}
\begin{frame}{Parité, inégalités et comparaison}
Vérifier le domaine avant la parité. Le signe de \(f\prime-g\prime\) renseigne sur la variation de \(f-g\), pas directement sur son signe.
\end{frame}
\begin{frame}{Parité, inégalités et comparaison}
Étudier la parité de \(x^4+x^2\) et \(x^3-x\) ; démontrer \(x^2+1\geq2x\).
\end{frame}
\begin{frame}{Optimisation}
Optimiser consiste à définir la variable et son domaine, exprimer la grandeur à optimiser, puis comparer toutes les valeurs candidates, y compris les bornes.
\end{frame}
\begin{frame}{Optimisation}
Une boîte a une base carrée de côté \(x\) et une hauteur \(6-x\) (la somme du côté et de la hauteur vaut 6), avec \(0\le x\le6\). Son volume \(V(x)=x^2(6-x)\) a pour dérivée \(V\prime(x)=12x-3x^2=3x(4-x)\), positive sur \([0;4]\) et négative sur \([4;6]\). Le volume maximal est \(V(4)=32\).
\end{frame}
\begin{frame}{Optimisation}
Une valeur numérique fournie par un algorithme n’est pas une preuve de minimum ni de convergence.
\end{frame}
\begin{frame}{Optimisation}
\begin{columns}[T,onlytextwidth]
\begin{column}{.57\textwidth}
Un rectangle de périmètre 12 m a largeur x. Déterminer l’aire maximale par dérivation.
\end{column}
\begin{column}{.4\textwidth}
\centering
\includegraphics[width=\linewidth,height=4.5cm,keepaspectratio]{figures/premiere-06-2.pdf}
\end{column}
\end{columns}
\end{frame}
\end{document}
